\chapter{黎曼电磁学}

本章把黎曼宇宙中"光"的概念发展为一套完整的电磁学理论。我们将看到，四维矢量势 $\bA$ 给出电磁场张量 $F$，电荷受力的规律与洛伦兹宇宙形式上完全平行，但偏振模式、多普勒频移、静电力的空间振荡乃至一切静电与磁现象都表现出本质性的差异。理解这些差异，是理解黎曼宇宙中物质如何相互作用的关键。

\section{三种偏振}

在我们所在的宇宙中，真空中的光是完全\textbf{横波}（transverse wave）：构成光波的振荡电场与磁场始终与该波在空间中传播的方向正交。

\orthofig{images/04__001.png}{横波与纵波}

在黎曼宇宙中，如矢量波一节的讨论，等价的光波在某种意义上也总是"横"的：描述振荡场的四维矢量 $\bA$ 与四维传播矢量 $\bk$ 正交。若波写作
\[
A(x) = A_0 \sin(\bk \cdot x),
\]
则必须有 $\bk \cdot \bA_0 = 0$。施加该条件后，波矢矢量仍允许三个独立方向。例如，设传播矢量 $\bk$ 位于 $xt$ 平面内（如右图）。那么当 $\bA_0$ 沿 $y$ 轴、沿 $z$ 轴、或位于 $xt$ 平面内且与 $\bk$ 成直角时，这三个方向互相垂直，且都与 $\bk$ 正交。（第二种可能性在图中未画出，因为为了显示 $t$ 方向我们省略了 $z$ 方向。）

然而把这些矢量投影到三维空间后，前两个方向\textbf{仍}保持与波传播方向正交，第三个方向却变得与它平行（或反平行）。在这个意义上，黎曼宇宙中"光"的等价物可以包含一个\textbf{纵模}（longitudinal mode）：场沿波传播的方向振荡。（图中我们展示的是向 $xy$ 平面的投影。由于所用矢量本就没有 $z$ 分量，其结果与投影到三维空间相同。）

必须强调，这三种\textbf{偏振}（polarisation）模式——两种看起来横、一种看起来纵——当我们把光本身作为四维空间中的平面波来看待时，是完全相同的。在每种情形下，$\bk$ 与 $\bA_0$ 都互相正交。只有当我们选定某个特定方向作为时间轴，并\emph{相对于该选择}去区分四维空间中的方向时，谈论"纵偏振"才有意义。尽管如此，这里与我们的宇宙仍存在非常真实的差别——我们的宇宙只有两种偏振。在黎曼宇宙中，如果偏振片的等价物移除了光的一个横向分量，那么"正交偏振片"（在我们的宇宙中会完全阻断光）仍会透射随机偏振光源中的纵分量。

\section{多普勒频移（Doppler Shift）}

\orthofig{images/04__024.png}{洛伦兹宇宙中的多普勒频移}

多普勒频移（Doppler shift）是一个众所周知的效应：观察者的运动改变了他们测得的光波波长与频率。

右图显示同一列波前扫过洛伦兹宇宙中三位不同的观察者，光从右侧射来。"参考观察者"假定相对光源静止，另外两位观察者则分别朝向、背离光源运动。

参考观察者测得相邻波前之间的周期为 OR。朝光运动的观察者在事件 O 和 F 处遇到波前，但图中所示的 OF 的欧几里得长度并非观察者经历的时间，后者需用洛伦兹版勾股定理计算。为便于比较，我们用了双曲线——洛伦兹宇宙中一条等固有时（proper time）曲线，正如圆弧是欧几里得空间中的等距曲线——把该间隔移至参考观察者的世界线上。在那里可以明确看到 $\mathrm{OF}'$ 小于 OR，故向前运动的观察者测得光的周期变短、频率升高。这称为\textbf{蓝移}（blue shift），因为蓝光处于可见光谱的高频端。

类似地，向后运动的观察者在事件 O 与 B 处遇到波前；为比较我们把该间隔移至 $\mathrm{OB}'$，此处它显然长于 OR。因此向后运动的观察者测得光的周期变长、频率降低。这称为\textbf{红移}（red shift）。

设 F 的坐标为 $(x_F, t_F)$。若观察者速度为 $v$，则 $x_F = v\,t_F$。记 $\tau_R$ 为参考观察者测得的周期，即 OR。由于波前从 F 到 R 的速度为 1（采用我们的单位），有
\[
\tau_R = t_F + x_F = t_F\,(1+v),
\]
\[
t_F = \frac{\tau_R}{1+v}.
\]
于是朝前运动的观察者测得的光周期由洛伦兹版勾股定理给出：
\[
\tau_F = \sqrt{t_F^{2} - x_F^{2}}
       = t_F\sqrt{1 - v^{2}}
       = \tau_R\frac{\sqrt{1 - v^{2}}}{1+v}
       = \tau_R\sqrt{\frac{1 - v}{1+v}}.
\]
洛伦兹蓝移的\emph{频率}之比恰为周期之比的倒数：
\[
\frac{\nu_F}{\nu_R} = \frac{\tau_R}{\tau_F} = \sqrt{\frac{1+v}{1-v}}.
\]
洛伦兹蓝移因子不依赖于光的具体频率。设速度 $v = 0.6$，即光速的百分之六十，则蓝移因子为 2，来自正前方的一切光相对静止时的频率都翻倍；同理，来自正后方的光频率减半。

黎曼情形稍复杂，但让我们仍从同样的情景开始：一位观察者接近光源，另一位背离光源——两者速度都低于光本身的速度。

\orthofig{images/04__025.png}{黎曼宇宙中的多普勒频移}

此情形下从图中可清晰看到，朝前运动的观察者测得的光周期 $\mathrm{OF}'$ \emph{大于} OR，而向后运动的观察者测得的周期 $\mathrm{OB}'$ \emph{小于} OR。

这是否意味着接近光的观察者看到"红移"？这里我们必须停下来追问：在黎曼宇宙中用何种语言描述可见光谱？这并非关于异宇宙生物能否体验相同的"红""蓝"感受的哲学疑虑——只要我们保持一致，尽可把他们的色觉随意翻译成我们的词汇。但无论黎曼宇宙居民对光的感觉如何，我们都不可能做出尊重光的一切客观物理属性的完美翻译。在我们的宇宙中，红光比蓝光既有更长的波长、也有更长的周期。在黎曼宇宙中，\emph{没有任何一种光}同时具有比另一种光更长的波长\emph{和}更长的周期。关系总是相反的：波长越长的光，周期反而越短。

因此我们必须作出选择。小说《Orthogonal》中采取的任意方案是，把可见光中最短的\emph{波长}译为"紫"，于是我们把"蓝移"一词定义为"朝向更短的波长"。

确定了术语之后，我们继续来量化最容易的东西，而这仍然是周期或频率之比而非波长之比。若该光相对参考观察者的速度为 $c$，则波前沿与传播矢量相反的方向以速度 $1/c$ 运动，因为与斜率为 $c$ 的直线垂直的直线的斜率是 $-1/c$。黎曼版的计算成为
\[
\tau_R = t_F - \frac{x_F}{1/c} = t_F\,(1 - vc),
\]
\[
t_F = \frac{\tau_R}{1 - vc},
\]
\[
\tau_F = \sqrt{t_F^{2} + x_F^{2}}
       = t_F\sqrt{1 + v^{2}}
       = \tau_R\frac{\sqrt{1 + v^{2}}}{1 - vc}.
\]

\orthofig{images/04__026.png}{多普勒频移后的频率}
\orthofig{images/04__028.png}{多普勒频移后的空间频率}

于是黎曼"蓝移"对接近光源的观察者而言，反而涉及观测频率的\emph{降低}：
\[
\frac{\nu_F}{\nu_R} = \frac{\tau_R}{\tau_F} = \frac{1 - vc}{\sqrt{1 + v^{2}}}.
\]
若观察者速度 $v$ 达到 $1/c$，则其测得的光频率降为零。图中世界线为 OG 的观察者就是例子：该观察者看到光速为无穷大，波长为最小可能值。

若观察者朝前运动的速度大于 $1/c$（例如世界线为 OH 的观察者），来自光源的传播矢量在他看来指向过去，因此在与光的任何相互作用中，他都会认为\emph{自己}是光源，传播矢量方向反转。

把同一公式中的 $v$ 取负值即可得到背离光源的观察者的频率比。从图中可见，当 $v = -c$ 时会发生特别的事：世界线为 OC 的观察者与光同步，使光的相对速度为零。对这样的观察者，频率达到最大值 $\nu_{\max}$——即可能的最大频率，因为他正径直穿过波前。

当观察者开始超过光速、比光更快地向后运动时又如何？我们的频率比公式给出一个正值，此时它小于最大值，故光应当仍可见。但当你以比光本身更快的速度逃离光源时，如何还能接收到来自光源的光？我们的假定是：光在一段延展的时间内被持续产生（比如星光就是如此），因此这不是观察者与某一单个光脉冲之间的赛跑。若观察者运动快于所讨论的光（如世界线为 OD 的观察者），他会开始追上"旧"光——那些已有足够领先量的光正等着他。并且他会看到这些光来自他行进的方向，而非光源本身的方向。

除了多普勒频移改变光的频率，几何关系还意味着运动的观察者测得的光线之间的\emph{夹角}与参考观察者不同，这一效应称为\textbf{光行差}（aberration）。在黎曼情形下，我们可用普通欧几里得几何理解多普勒频移与光行差。

设参考观察者看到天空中均匀分布的大量光源发出某一特定颜色的光。所有入射光线在四维空间中构成一个锥面，其轴为参考观察者的世界线。若把所有传播矢量投影到参考观察者所认为的空间中，它们会对称地展开成一个圆（下图中灰色所示）。

位于同一位置、以某速度 $v$ 向右运动的观察者，看到的光来向对应于同一锥面上这些光线在他们自己空间中的投影——该空间相对参考观察者的空间是倾斜的。这会产生两个效应：改变传播矢量投影的长度，并改变这些投影之间的角度。传播矢量在观察者空间上投影的\emph{长度}正比于光的空间频率（如最初考察波时所讨论的），因此更长的投影意味着更短的波长。

\orthofig[0.72\textwidth]{images/04__027.png}{黎曼宇宙中的多普勒频移与光行差}

该图展示了观察者取各种速度时的结果，假定所有入射光在参考观察者参考系中的速度为 0.75。图中颜色与小说用于可见色调的翻译方案一致，虚线表示紫外与红外。

第一个例子取 $v = 0.25$：来自观察者行进方向的光发生蓝移，光源在天空中显得被压缩挤在一起；来自相反方向的光发生红移，在天空中散开。用这种说法——把红移、蓝移理解为波长而非频率——这段描述听起来与洛伦兹版本一模一样。

一旦 $v = c$，一切光看起来都来自前方。当 $v = 1/c = 4/3$（未画出）时，正前方的光线将逆着时间到达，因而不可见；当 $v = 2$ 时，锥面的可觉察部分因此被略去。

对 $v = 2$，我们看到一件在洛伦兹物理中绝不可能发生的奇事：来自同一表观方向、具有两种不同颜色、实际上来自两个不同光源的光。在红、黄、绿光线与紫外光混杂之处，较红的光实际上来自后方（按参考观察者而言），但运动观察者的倾斜使其看起来来自前方，并与真正来自前方的光叠加在一起。

当 $v = \infty$ 时，参考观察者认为来自前方的一切光如今都来自运动观察者的未来，因而不可见；此时可见的光全部来自后方，并被挤入一片天空区域，其角宽度恰好等于入射光锥的角宽度。

请记住，我们这里描述的一切都只涉及在参考观察者看来\emph{只有一种颜色}的光。即便对一位相对恒星平均运动静止的观察者而言，在黎曼宇宙中大多数恒星也会呈现为\emph{色彩拖尾}（colour trails），因为恒星在发出目前到达观察者的红光之时与稍后发出同样正在到达的、更快的紫光之时之间，已经移动了一点。多普勒频移与光行差如何改变这些恒星拖尾，在小说第一卷中有所讨论；书中预告片可见这些效应的图示，另有一个相关主题的视频教程。

\begin{notebox}[术语约定]
黎曼宇宙中波长与周期的关系与我们的宇宙相反，无法同时把"更长的波长"和"更长的周期"都译为"更红"。小说采取任意方案：把最短波长译为"紫"，并约定\textbf{蓝移＝朝向更短的波长}。
\end{notebox}

\section{黎曼电磁场}

至此我们一直在谈论一个"振荡场" $\bA$，却未确切说明它是什么。事实证明，可以完全类比我们自己宇宙的洛伦兹电磁学，发展出一整套黎曼电磁学理论，其中包括我们可以合理地称为电场与磁场的矢量场，以及一种类似电荷的物质属性——它使相关物质受到电磁力，并具备产生电磁波的能力。

在这一框架中，我们一直称为 $\bA$ 的场就是所谓的\textbf{四维矢量势}（four-vector potential）。正如电势是这样一种量：它在空间中的变化率给出电场；四维矢量势则是一个矢量，其导数给出完整的电磁场。具体地说，电磁场 $F$ 的分量为
\[
F_{ij} = \partial_i A_j - \partial_j A_i,
\]
其中 $i$、$j$ 取遍 $x$、$y$、$z$ 与 $t$。换言之，$F$ 由一个 $4\times4$ 矩阵描述，它由四维矢量势 $\bA$ 导出。

这个矩阵的意义是什么？设有带电量为 $q$ 的粒子，沿一条世界线运动，其四维切线是单位矢量 $\bu$，即粒子的\emph{四维速度}（four-velocity）。若把 $\bu$ 乘以电荷 $q$ 与矩阵 $F$，就得到带电粒子所受的力 $\mathbf{f}$：
\[
\mathbf{f} = q F \bu.
\]
当然，在牛顿力学中力是三维矢量，但在相对论物理中（无论黎曼还是洛伦兹），把力视为四维矢量更为合理。正如牛顿的著名方程
\[
\mathbf{f} = m\,\mathbf{a}
\]
（力等于质量乘加速度），力的四维矢量 $\mathbf{f}$ 与加速度的四维矢量 $\mathbf{a}$ 相联系……后者又与四维速度通过
\[
\mathbf{a} = \partial_{\tau}\bu
\]
相联系。四维加速度是四维速度相对\emph{固有时}（proper time）$\tau$ 的变化率，$\tau$ 是沿粒子世界线测得的时间。把所有这些片段放在一起，我们得到关于电磁场 $F$ 对带电量 $q$、质量 $m$ 的粒子运动之影响的如下描述：
\[
\partial_{\tau}\bu = \frac{q}{m} F \bu.
\]
四维速度 $\bu$ 的模长恒为 1。任何与 $\bu$ 自身平行的改变都会改变其模长，因此描述其变化率的矢量必须始终与 $\bu$ 正交。这意味着 $F$ 必须属于一类特殊的矩阵：用 $F$ 乘任何矢量，所得结果都与原矢量正交。不难看出，这对任何\emph{反对称矩阵}（antisymmetric matrix）都成立——即等于自身转置之负的矩阵，或 $F_{ij} = -F_{ji}$。于是，使用爱因斯坦求和约定，有
\[
\bu \cdot (F\bu) = u^{i} F_{ij} u^{j} = u^{i} (-F_{ji}) u^{j} = -\bu \cdot (F\bu),
\]
这表明点积必须为零。而用 $\bA$ 定义 $F$ 的方式，保证了它\emph{确实}是反对称的。

那么，带电粒子在\emph{恒定}电磁场中运动时将走怎样的世界线？最容易考虑的情形是矩阵 $F$ 取特别简单的形式：
\[
F_{ij} = b_i c_j - b_j c_i,
\]
其中 $\mathbf{b}$ 与 $\mathbf{c}$ 是任意两个\emph{不平行}的矢量。（若 $\mathbf{b}$ 是 $\mathbf{c}$ 的倍数，则右端变为零。）我们把此式缩写为
\[
F = \mathbf{b} \wedge \mathbf{c}.
\]
符号"$\wedge$"称为"楔积"（wedge product）。此时可把 $F$ 与矢量 $\bu$ 的矩阵乘积用 $\mathbf{b}$ 与 $\mathbf{c}$ 重写：
\[
F\bu = (\mathbf{c}\cdot\bu)\,\mathbf{b} - (\mathbf{b}\cdot\bu)\,\mathbf{c}.
\]
显然这是始终位于 $\mathbf{b}$、$\mathbf{c}$ 张成平面内的矢量；它也只依赖于 $\bu$ 落在该平面内的分量：若 $\bu$ 完全与该平面正交（即 $\mathbf{b}\cdot\bu = \mathbf{c}\cdot\bu = 0$），则 $F\bu$ 为零。

\orthofig{images/04__002.png}{恒定电磁场中的粒子}

结果一般形如右图。当 $\bu$ 既不完全落在 $\mathbf{b}$、$\mathbf{c}$ 张成的平面内，又不完全与之正交时，粒子会以恒定速率保持在与该平面正交的方向上运动，而它在平行于该平面的方向上的运动则被弯成一个圆。于是粒子世界线整体是一条螺旋线。

若该平面由你视为空间的两个方向张成（换言之，平面与你的世界线正交），那么运动的圆周部分在你看来就是空间中的运动：粒子要么在一个固定平面内打转，要么在空间中沿螺旋前进。这正是我们的宇宙中与\textbf{磁场}相联系的行为。事实上，在一种略为老式的粒子加速器——回旋加速器（cyclotron）中，正是用恒定磁场使粒子如此打转。磁场的描述约定是：赋予它一个与粒子打转平面正交的空间方向。

若该平面包含你视为时间的方向，则世界线的弯曲将表现为粒子在平面内另一（空间）方向上加速。在我们宇宙中，带电粒子被如此加速时，我们认为这是\textbf{电场}的作用。在洛伦兹物理中，加速粒子的世界线是双曲线而非圆，但除此之外效应相同。我们赋予电场一个落在该平面内的空间方向。

于是形如 $F = \mathbf{b}\wedge\mathbf{c}$ 的电磁场被视为：
\begin{itemize}
\item 若 $\mathbf{b}$ 与 $\mathbf{c}$ 都是空间方向，即一个同时垂直于 $\mathbf{b}$、$\mathbf{c}$ 的磁场；
\item 若该平面包含你的时间方向，即一个空间方向落在 $\mathbf{b}$、$\mathbf{c}$ 张成平面内的电场。
\end{itemize}
若该平面不干净地落入这两类中的任何一类，$F$ 将被视为同时包含电场与磁场。

让我们用这些描述来看黎曼光的横模与纵模中发生了什么。一般地，若
\[
A(x) = A_0 \sin(\bk\cdot x),
\]
则电磁场由矩阵 $F$ 描述：
\[
F(x) = (\bk \wedge A_0)\cos(\bk\cdot x).
\]
设光传播矢量位于 $xt$ 平面内，在 $x$ 方向速度为 1，即 $\bk = \mathbf{e}_x + \mathbf{e}_t$。（其实它必须是满足条件 $|\bk| = \wms$ 的该矢量的某一倍数，但 $\wms$ 的取值在下面并不重要。）

首先考察一个我们视为横模的波，取 $\bA_0 = \mathbf{e}_y$。$\bk$ 与 $A_0$ 张成的平面既不完全空间化，也不包含我们的时间轴 $\mathbf{e}_t$，故该波将\emph{同时}包含电场与磁场。但我们可以把 $F$ 拆成在我们的坐标系中分别对应磁场与电场的两部分：
\[
\begin{aligned}
F(x) &= \big((\mathbf{e}_x + \mathbf{e}_t)\wedge \mathbf{e}_y\big)\cos(\bk\cdot x) \\
     &= (\mathbf{e}_x \wedge \mathbf{e}_y)\cos(\bk\cdot x) + (\mathbf{e}_t \wedge \mathbf{e}_y)\cos(\bk\cdot x).
\end{aligned}
\]
第一部分，平面由 $\mathbf{e}_x$、$\mathbf{e}_y$ 张成，是一个指向 $z$ 方向的振荡磁场。第二部分，平面由 $\mathbf{e}_t$、$\mathbf{e}_y$ 张成，是一个指向 $y$ 方向的振荡电场。这与我们自己宇宙中的光非常相似：电场、磁场与光的传播方向三者两两垂直。

其次考察一个我们视为纵模的波，取 $\bA_0 = \mathbf{e}_x - \mathbf{e}_t$。
\[
\begin{aligned}
F(x) &= \big((\mathbf{e}_x + \mathbf{e}_t)\wedge(\mathbf{e}_x - \mathbf{e}_t)\big)\cos(\bk\cdot x) \\
     &= 2\,\mathbf{e}_t \wedge \mathbf{e}_x \cos(\bk\cdot x).
\end{aligned}
\]
在化简楔积时，我们利用了任何矢量与自身的楔积为零、以及 $\mathbf{b}\wedge\mathbf{c} = -\mathbf{c}\wedge\mathbf{b}$ 这两点。结果是沿 $x$ 方向——即平行于光本身运动的方向——的一个振荡的\emph{纯电场}。

\begin{keybox}[黎曼电磁场的要点]
四维矢量势 $\bA$ 通过 $F_{ij} = \partial_i A_j - \partial_j A_i$ 给出反对称的电磁场张量 $F$；其反对称性保证四维速度模长恒定。带电荷 $q$、质量 $m$ 的粒子的运动满足 $\partial_\tau \bu = (q/m)F\bu$。当 $F$ 的平面由两个空间方向张成时是磁场，当包含时间方向时是电场，否则两者兼有。
\end{keybox}

\section{波纹状的库仑力（The Corrugated Coulomb Force）}

黎曼矢量波方程（RVW）支配真空中电磁场的四维矢量势 $\bA$——即不存在任何可能产生该场的带电物质时的情形。但我们可以轻易地把 RVW 方程推广，加入一个描述电荷存在的"源"项，如下所示：
\begin{equation}
\partial_x^{2}\bA + \partial_y^{2}\bA + \partial_z^{2}\bA + \partial_t^{2}\bA + \wms^{2} \bA + \bj = \mathbf{0}. \tag{RVWS}
\end{equation}

我们在此方程中加入的四维矢量 $\bj$ 有非常简单的含义。设有某种带电物质存在于所关注的区域，一位相对该物质静止的观察者测得其电荷密度为 $\rho$。若该物质的四维速度为 $\bu$，我们定义
\[
\bj = \rho\,\bu.
\]
现在，$\bu$ 是切于物体世界线的单位矢量。但在世界线上任一点都有\emph{两个}这样的矢量，指向相反的方向。我们如何选取正确的那一个？若某观察者以其中一种选择作为其时间轴，带电物质会显得带正电；若取另一方向，同一物质会显得带负电。改变选择会把 $\rho$ 与 $\bu$ 同乘以 $-1$，从而抵消这一变化。因此选哪一个都无所谓。

我们把矢量 $\bj$ 称为\textbf{四维电流}（four-current）。若一位并不相对带电物质静止的观察者测量 $\bj$ 的分量，他将同时看到一个电荷密度 $j^t$，以及一个描述电荷运动的\textbf{电流密度}（current density）$j^x$、$j^y$、$j^z$。

该方程最简单也最重要的解之一称为\textbf{库仑势}（Coulomb potential）。这是带电量为 $q$ 的粒子占据单个点、并在某坐标系中静止时的势。当然这是一种理想化，它有沿粒子世界线出现无穷大电荷密度这样一个略显麻烦的数学特征。尽管如此，求解并不太难。如何做到这一点的细节见该主题的补充注记；这里我们把它作为已知结果接受。若 $r$ 为到粒子的距离，则四维矢量势为
\[
\bA(r) = \frac{q}{4\pi r}\cos(\wms r)\,\mathbf{e}_t.
\]
当 $\bA$ 只有这样的一个 $t$ 分量、且 $\bA$ 自身不随时间变化时，更简单的做法是忘掉四维矢量势，只用电势 $\phi$ 来工作，它是 $A_t$ 的相反数。

\emph{库仑势：}
\[
\phi(r) = \begin{cases}
-\dfrac{q}{4\pi r}\cos(\wms r), & \text{（Riemannian）}\\[8pt]
\dfrac{q}{4\pi r}, & \text{（Lorentzian）}
\end{cases}
\]
电势是带一单位电荷的正"试探粒子"的势能——正如任何随处变化的势能一样，它产生一个力。力在任一方向的分量等于 $\phi$ 在该方向变化率的相反数（粒子将朝势能\emph{增加}的相反方向加速，正如石头会沿山\emph{滑下}）：
\[
\mathbf{f} = -\big(\partial_x\phi\,\mathbf{e}_x + \partial_y\phi\,\mathbf{e}_y + \partial_z\phi\,\mathbf{e}_z\big),
\]
用梯度符号"$\nabla$"缩写为
\[
\mathbf{f} = -\nabla\phi.
\]
我们声称的 $\phi$ 与 $A_t$ 之间的简单关系，可通过检验二者对试探粒子给出完全相同的力来验证。

\orthofig{images/04__003.png}{正电荷周围黎曼与洛伦兹电势的比较}
\orthofig{images/04__004.png}{负电荷周围黎曼与洛伦兹电势的比较}

$\phi$ 是什么样子？右上图显示在黎曼与洛伦兹电磁学中正电荷周围电势的情形，下图显示负电荷周围的电势。

黎曼情形与洛伦兹情形之间有两个显著差异。在我们自己的宇宙中，众所周知\emph{异种电荷相吸}，\emph{同种电荷相斥}。但在黎曼宇宙中，正电荷附近的电势是一个井而非一座山，因此若两个正电荷足够靠近，它们会相吸；异种电荷若足够靠近则会相斥。

第二个差异是，我们必须指明两电荷之间的距离，才能说它们相吸还是相斥。在洛伦兹宇宙中，静电力只是随距离逐渐减弱；在黎曼宇宙中，力的方向还随距离\emph{交替}变化。从 $\phi$ 解的形式可以看出，这种交替的空间频率等于光的最大可能频率 $\nu_{\max}$（因为 $\wms = 2\pi\nu_{\max}$）。等价地，其起伏的波长等于光的最小可能波长 $\lambda_{\min} = 1/\nu_{\max}$。

不难验证我们对 $\bA$ 的解满足 RVWS 方程；这种交替电势其实不大令人惊讶，因为我们已经看到，紫外极限下的平面波也涉及一个不随时间变化、而以频率 $\nu_{\max}$ 在空间中振荡的电磁场。该方程要求所有方向上频率的平方和为常数，因此任何不随时间变化的静场\emph{必然}在空间中振荡。

但我们怎能确信方程本身正确？若把方程中 $\bj$ 的符号反号，我们就会得到原解的相反情形，从而得到近距相斥的同种电荷。是什么阻止我们提出黎曼宇宙中电磁学定律取\emph{那种}形式呢？

最简单的回答是回想我们在能量与动量注记中讨论过的一件事：在黎曼宇宙中相对论总能量的意义与势能和动能相反。在洛伦兹宇宙中，电磁场拥有能量，正比于场强的平方。给定一个正电荷与一个负电荷，当二者相互靠近时该能量会减少，因为它们会越来越多地\emph{抵消}彼此的场。但在黎曼宇宙中，使势能最小意味着使场的能量\emph{最大}。对于足够靠近的同种电荷，把它们拉得更近会使它们形状相同的场叠加得越来越匹配，\emph{减少}两粒子场之间的抵消，从而\emph{增加}整体场能。

\section{不受欢迎的电子学（Unpopular Electronics）}

在我们的宇宙中，环绕带电粒子（如电子）的微小电场能够轻易累积到产生可见效应的程度。把塑料笔在合适的衣物上摩擦十几次，它就会带上多余电荷，产生强到足以克服重力吸起一小片纸屑的电场。（纸张本身并无净电荷，但笔的电场使纸表面的电荷分离。）多余电荷在笔表面如何精确分布无关紧要；只要\emph{存在}多余电荷，其每一部分都会协同增强笔周围的电场强度——而这个电场在时间上持久、在空间上一致，在很长的区域内大致指向同一方向，并保持到电荷缓慢泄漏为止。

\orthofig{images/04__005.png}{随机排列电荷产生的电势}
\orthofig{images/04__006.png}{晶格排列电荷产生的电势}

在黎曼宇宙中，由于库仑势在空间中的振荡，多余电荷的效应非常不同。大量相同粒子杂乱排列（如右图）时，对附近某点电势的贡献既有正也有负，彼此大量抵消。当然抵消不会完美，但残余的场将是一系列密布的峰谷，而非洛伦兹情形下那种单一、持久的斜坡。

有序排列的电荷阵列（如右下图）在某些点可能比随机排列抵消得更少，但同样会产生起伏的电势地形。正试探粒子若从正确的初始位置静止出发，确实可能被这堆电荷排斥，但该效应不像洛伦兹物理中的等价物那样容易复现。更有甚者，任何置于该场中、由许多粒子（而非单一粒子）组成的物体，将经历在诸多略有不同位置上取决于场的各个力的总和，从而导致进一步的抵消。

其结果是：静电现象将对电荷的细节、微观排列十分敏感，且总体而言远不如在我们宇宙中那样容易探测和利用。

类似的效应也适用于磁性。在我们的世界里，仅需把大量小磁偶极子的方向对齐，它们的场便足以互相增强；永磁体正是这样工作的。但在黎曼宇宙中，即便偶极子本身全部对齐，当你穿越空间时，每个偶极子产生的场也会在两个相反方向之间来回切换，且在任意一点，来自距离略有不同的各个偶极子的贡献之间都会发生抵消。

不过，尽管这些问题排除了那种让我们最初窥见电磁学本质的、容易的低技术实验，它们并非完全不可逾越。我们描述过的静场（不随时间变化）其波长极小，正是这一点使一切如此精细、如此容易抵消；而一旦引入\emph{时间}上的振荡，这种小波长会逐渐让位于更大的波长。在直流电路中需要极高精度才能工作的电子元件，若改用频率足够高的\textbf{交流电}（alternating current），对不完美的容忍度将大得多。

什么叫"频率足够高"？只有当频率 $\omega$ 达到 $0.995\,\wms$ 时，场的波长才会增长到最小波长的十倍，因此要获得任何可用的空间不精确容忍度，代价就是必须处理黎曼宇宙中能传播的、本质上频率最高的电磁辐射。这将带来它自己的技术挑战——但它们与在我们宇宙中处理高频的困难并不完全相同，后者有些实际上与波长有关。

对不介意稍高数学水平的读者，这个主题另有补充材料。
